
This appendix details the operational interface elements and chart rendering structures implemented on the interactive web dashboard.

\section{D.1 Anomaly Mapping Canvas}
\begin{itemize}
  \item \textbf{Basemap Tile Layer}: CartoDB Dark Matter tiles (\texttt{https://{s}.basemaps.cartocdn.com/dark\_all/{z}/{x}/{y}{r}.png}) are chosen to reduce visual noise and highlight fluorescent marker indicators.
  \item \textbf{Pulsing Markers}: Custom Leaflet HTML \texttt{divIcon} elements are styled using CSS keyframe expansions. The pulsing ring expands dynamically:
    \texttt{}`css
    @keyframes markerPulse {
        0\% { transform: scale(0.5); opacity: 1; }
        100\% { transform: scale(2.5); opacity: 0; }
    }
    \texttt{}`

\end{itemize}
\section{D.2 Dual-Canvas Image Split Slider}
\begin{itemize}
  \item To render comparison tiles without loading heavy packages, two HTML5 \texttt{<canvas>} elements are loaded.
  \item The baseline image is painted to the lower canvas, and the anomaly image is painted to the upper canvas.
  \item An absolute-positioned divider handle listens to horizontal pointer drag events (\texttt{mousemove} / \texttt{touchmove}) to calculate split percentages.
  \item The anomaly canvas is dynamically cropped using the CSS \texttt{clip-path: inset(0 0 0 splitXpx)} rule, revealing the baseline on the left half and the anomaly on the right half.

\end{itemize}
\section{D.3 Spectral Profile Radar Chart}
\begin{itemize}
  \item Rendered using Chart.js inside a responsive glass-container.
  \item Maps four axis points representing NDVI, NDWI, BSI, and SWIR reflectance.
  \item The selected pixel's signature is drawn in fluorescent cyan, contrasted with the baseline mean drawn in deep blue, allowing immediate visual verification of land-clearing vs. water-flooding anomalies.

\end{itemize}